% catalog.tex — Section 4: Model Catalog
% Included by main.tex as % catalog.tex — Section 4: Model Catalog
% Included by main.tex as % catalog.tex — Section 4: Model Catalog
% Included by main.tex as % catalog.tex — Section 4: Model Catalog
% Included by main.tex as \input{sections/catalog}
% No \begin{document} / \end{document}

\subsection{Task Types}
\label{sec:catalog:tasks}

VisionServe organises its model portfolio around seven computer-vision task types.
Table~\ref{tab:task-types} defines each type, states the expected input and output
contract, and names representative models currently registered in the system.

\begin{table*}[t]
  \centering
  \small
  \caption{Task types supported by VisionServe.}
  \label{tab:task-types}
  \begin{tabular}{llp{6.2cm}l}
    \toprule
    \textbf{Task} & \textbf{I/O Contract} & \textbf{Description} & \textbf{Example Models} \\
    \midrule
    \texttt{detection}     & image $\to$ boxes     &
      Localise and classify objects in an image.
      Returns axis-aligned bounding boxes in original-image coordinates
      together with class labels and confidence scores.
      &
      RF-DETR, SCRFD, PaddleOCR \\[4pt]

    \texttt{segmentation}  & image + prompt $\to$ masks &
      Produce per-pixel binary masks, optionally conditioned on a box or
      point prompt.  Masks are encoded as column-major run-length encoding
      (RLE) in the unified \texttt{Result} schema.
      &
      MobileSAM, EfficientSAM, SAM2 \\[4pt]

    \texttt{open\_vocab}   & image + text $\to$ boxes &
      Detect objects whose categories are specified as free-form natural-language
      queries at request time, without any fixed class vocabulary.
      &
      GroundingDINO \\[4pt]

    \texttt{classification}& image $\to$ label + score &
      Assign a single top-\(k\) label from a closed-set vocabulary
      (e.g.\ ImageNet-1k) to the entire image.
      &
      EfficientNet-B0, MobileNetV3 \\[4pt]

    \texttt{depth}         & image $\to$ depth map &
      Estimate a per-pixel relative or metric depth map.  Output is a
      single-channel float32 tensor aligned to the input resolution.
      &
      MiDaS, Depth Anything V2 \\[4pt]

    \texttt{embed}         & image $\to$ vector &
      Extract a compact, semantically rich embedding vector suitable for
      retrieval, zero-shot classification, or downstream ML pipelines.
      &
      CLIP ViT-B/32 \\[4pt]

    \texttt{grasp}         & image $\to$ grasps &
      Synthesise planar parallel-jaw grasps ($(x,y,\theta,\text{width},\text{quality})$ in
      original-image coordinates) from a segmentation mask via a pure-Go analytic
      antipodal force-closure search (no extra ONNX weights).  Class-agnostic
      (whole-image automask) or class-aware (detector $\to$ mask $\to$ grasp).
      &
      grasp, grasp-rfdetr, grasp-gd \\
    \bottomrule
  \end{tabular}
\end{table*}

\subsection{Model Inventory}
\label{sec:catalog:inventory}

Table~\ref{tab:model-inventory} enumerates all 19 models in the current VisionServe
registry.  Each model is distributed as one or more ONNX files; multi-session models
(MobileSAM, EfficientSAM, SAM2, Grounded-SAM, and the grasp pipelines) use the
\texttt{PipelineModel} interface and declare a \texttt{files:} role map in their manifest.
The three grasp pipelines reuse the segmentation/detection weights of sibling models and add
no new ONNX file: the grasp synthesis itself is a pure-Go analytic stage.

\begin{table*}[t]
  \centering
  \small
  \caption{%
    VisionServe model inventory.
    Architecture notes describe the primary network family.
    Models marked $\dagger$ are pending benchmark runs (see text).%
  }
  \label{tab:model-inventory}
  \begin{tabular}{llll}
    \toprule
    \textbf{Model} & \textbf{Task} & \textbf{License} & \textbf{Architecture} \\
    \midrule
    \texttt{rf-detr}          & detection     & Apache-2.0 & ResNet-50 backbone + DETR head (NMS-free) \\
    \texttt{rf-detr-nano}     & detection     & Apache-2.0 & Lightweight DETR variant \\
    \texttt{rt-detr}$^\dagger$& detection     & Apache-2.0 & Hybrid CNN/Transformer real-time DETR \\
    \texttt{scrfd}            & detection     & MIT        & Sample- \& computation-redistributed face detector \\
    \texttt{paddle-ocr}       & detection     & Apache-2.0 & PP-OCRv4 text detection backbone \\
    \texttt{grounding-dino}   & open\_vocab   & Apache-2.0 & Swin-T image encoder + BERT text encoder \\
    \texttt{mobile-sam}       & segmentation  & Apache-2.0 & Tiny ViT encoder + SAM decoder \\
    \texttt{efficient-sam}    & segmentation  & Apache-2.0 & ViT-S/16 encoder with masked-image pretraining \\
    \texttt{sam2}             & segmentation  & Apache-2.0 & Hiera-Tiny image/video encoder + SAM2 decoder \\
    \texttt{nano-sam}$^\dagger$& segmentation & Apache-2.0 & ResNet-18 distilled SAM encoder \\
    \texttt{grounded-sam}    & open\_vocab    & Apache-2.0 & GroundingDINO boxes $\to$ MobileSAM masks (composite) \\
    \texttt{midas}            & depth         & MIT        & DPT-Hybrid (ViT-H + dense prediction transformer) \\
    \texttt{depth-anything-v2}$^\dagger$ & depth & Apache-2.0 & ViT-S teacher-student depth foundation model \\
    \texttt{efficientnet-b0}  & classification& Apache-2.0 & EfficientNet-B0 compound-scaled CNN \\
    \texttt{mobilenet-v3}     & classification& Apache-2.0 & MobileNetV3-Large depthwise-separable CNN \\
    \texttt{clip}             & embed         & MIT        & CLIP ViT-B/32 vision encoder (image side only) \\
    \texttt{grasp}            & grasp         & Apache-2.0 & MobileSAM automask $\to$ pure-Go analytic grasps (class-agnostic) \\
    \texttt{grasp-rfdetr}     & grasp         & Apache-2.0 & RF-DETR boxes $\to$ MobileSAM $\to$ analytic grasps (class-aware) \\
    \texttt{grasp-gd}         & grasp         & Apache-2.0 & GroundingDINO boxes $\to$ MobileSAM $\to$ analytic grasps (text-prompted) \\
    \bottomrule
  \end{tabular}
\end{table*}

\noindent
Three models marked $\dagger$ remain pending: \texttt{rt-detr} and
\texttt{depth-anything-v2} require authenticated HuggingFace downloads, and
\texttt{nano-sam} requires a manual ONNX export.  All four grasp/composite pipelines
(\texttt{grasp}, \texttt{grasp-rfdetr}, \texttt{grounded-sam}, \texttt{grasp-gd}) are now
benchmarked (Section~\ref{sec:eval:latency}); the two GroundingDINO-based pipelines are
serialized for concurrency-safety (a fixed \texttt{SA\_ONSTACK} crash, Section~\ref{sec:limitations})
and are reported single-stream.

\paragraph{Grasp and license discipline.}
The grasp capability also illustrates why the license gate (C1) is load-bearing rather than
decorative.  In selecting a grasp method we rejected the most accurate learned grasp
models, GQ-CNN / FC-GQ-CNN (Dex-Net), under a UC~Berkeley ``educational, research, and
not-for-profit'' license, because that non-commercial bar is incompatible with
VisionServe's permissive, commercial-friendly stance, as AGPL is.  The shipped path
instead pairs a permissive MobileSAM mask with a clean-room pure-Go analytic synthesizer (no
weights, no license at all), so the seventh task adds capability without adding license risk.

\subsection{License Safety}
\label{sec:catalog:license}

Open-source computer-vision models carry heterogeneous licenses, and the distinction
between permissive and copyleft licenses has material consequences for downstream
deployments.  The most prominent example is the Ultralytics YOLO family
\cite{ultralytics2023yolo}, which is released under the GNU Affero General
Public License v3.0 (AGPL-3.0) \cite{agpllicense}.  AGPL-3.0 extends the GPL copyleft
to networked services: any software that uses AGPL code and provides access to
users over a network must release its entire source under AGPL-3.0.  As a result,
integrating a single AGPL model into an inference server would virally relicense the
server (and every product built on it) under AGPL-3.0.  This is incompatible with
closed-source, proprietary, or commercial deployments; the only escape is a separate
commercial license negotiated with the copyright holder.

VisionServe enforces an allowlist of permissive licenses
(\texttt{Apache-2.0}, \texttt{MIT}, \texttt{BSD-2-Clause}, \texttt{BSD-3-Clause})
in the registry parser at process startup \cite{apachelicense}.  Any model manifest
that declares a \texttt{license} field outside this allowlist causes the server to
abort with a clear error message, preventing accidental integration of copyleft models.
This check is implemented in code, not merely documented: the manifest
\texttt{license} field is a required string, and the parser rejects absent or
non-allowlisted values before any ONNX session is initialised.  This design choice is
deliberate: because VisionServe is itself Apache-2.0 and targets edge deployments in
commercial and closed-source products, its permissive character must be preserved by
every model it ships.

\paragraph{Threat model.}
We scope the license gate (contribution~C1), since its value depends on
stating both what it defends against and what it cannot.
\emph{In scope.}  The gate is designed to prevent four failure modes at the moment the
runtime loads model weights:
(i)~accidentally loading a copyleft (AGPL) model that would virally relicense an otherwise
permissive deployment;
(ii)~``permissive-washing'' or license drift, where a model's restrictive license is
silently dropped or relabelled as permissive on the way into the deployment;
(iii)~tampered or swapped weight bytes that no longer correspond to the audited model
(a content mismatch);
(iv)~weights pulled from an unvetted origin; and
(v)~a manifest relabelled by its author with a permissive-but-wrong license that diverges from
the maintainer-audited upstream (caught by the provenance ledger, which the allowlist alone
cannot detect because the wrong label is itself permissive).
\emph{Mechanism.}  Enforcement is entirely offline, in a single binary, and consists of
four bound checks applied at load time.
First, the declared \texttt{license} field must be in the permissive allowlist
(\texttt{Apache-2.0}/\texttt{MIT}/\texttt{BSD}); the registry parser
(\texttt{Manifest.validate()}) refuses any other value before a session is created.
Second, the SHA-256 digest of the weight bytes must equal the digest declared in the
manifest (\texttt{Manifest.VerifyWeights()}); a mismatch refuses the load.
Third, when a verified-source allowlist is supplied, the manifest's \texttt{source\_url}
must fall under an audited prefix.
Fourth, the check that closes the upstream-relabelling hole for the curated
catalog: in the opt-in verified mode the contributor-declared license must equal the
license recorded for that \texttt{source\_url} prefix in a \emph{maintainer-audited
license-provenance ledger} (\texttt{Manifest.VerifyLicenseProvenance()}).  The ledger is a
separate in-binary table in which a maintainer records, per audited upstream, the SPDX license
they verified by reading the upstream's actual LICENSE file, the URL of that file, and its
SHA-256 as evidence.  Because contributors edit manifests and maintainers edit the ledger, the
two authorities are separated: a pull request that flips a manifest's \texttt{license} to
sneak in a copyleft model is refused by a record the PR author does not control.  Together these
checks bind a declared-permissive license to specific audited bytes, an audited origin,
and an independent audited record of the upstream's true license, without any network call.
\emph{Out of scope.}  The allowlist and SHA-256/source pins alone still trust the declared
license string and mitigate only substitution (swapped bytes, wrong mirror), not
upstream mislabelling; the ledger removes that residue for the shipped catalog by
cross-checking against an independent human audit.  What remains out of scope is genuine: the
ledger's ground truth is a one-time human audit of the upstream LICENSE file (a re-audit detects
later divergence via the pinned LICENSE-file hash, but an upstream mislabelled at audit
time is not caught), and it covers only audited upstreams, not arbitrary un-audited uploads, for
which a universal license oracle is impossible.  We therefore frame~C1 as load-time
license-policy enforcement, hardened by content-hashing, a verified-source pin, and an
audited provenance ledger: strong for the curated catalog, and not a
machine-checkable license guarantee for arbitrary models.
\emph{Why existing layers do not cover it.}  As discussed in
Section~\ref{sec:related}, build-time software-composition-analysis scanners
(ScanCode, FOSSology, Snyk) operate on source dependency trees and emit reports, not a
load-time decision over weights; deploy-time admission controllers (Kyverno, OPA
Gatekeeper, Sigstore policy-controller) verify container-image signatures at pod
admission, i.e.\ container integrity, not a license policy over model weights; and model
signing (OpenSSF/Sigstore) proves signer identity and byte integrity, not license policy.
None of these enforces a license policy at the moment the runtime loads model weights,
offline, which is the gap~C1 fills.

\paragraph{Adversarial demonstration.}
To show that the gate refuses bad models rather than merely documenting an
intent to, Table~\ref{tab:gate} reports a reproducible adversarial demonstration.
Starting from a baseline manifest (permissive license, correct SHA-256, audited source)
that is admitted, we apply one adversarial mutation at a time and record the gate's
verdict.  Each mutation maps to one of the four bound checks above, and each is refused
with the runtime's own diagnostic message.  The six cases are exercised by the test
\texttt{TestGate\_AdversarialDemo}; the transcript is reproducible via
\texttt{go test ./internal/registry -run TestGate\_AdversarialDemo -v}.

\begin{table}[t]
  \centering
  \footnotesize
  \setlength{\tabcolsep}{3pt}
  \caption{%
    Reproducible adversarial demonstration of the load-time license gate (C1).
    Starting from an admitted baseline, each row applies one manifest mutation and
    records the runtime's verdict and (abbreviated) diagnostic.  Reproducible via
    \texttt{go test ./internal/registry -run TestGate\_AdversarialDemo -v}
    (\texttt{internal/registry/gate\_demo\_test.go}); the long digest is abbreviated
    with \(\dots\) and the temporary weight path is shown as \texttt{<model.onnx>}.%
  }
  \label{tab:gate}
  \begin{tabular}{p{1.2cm}p{1.5cm}p{5.4cm}}
    \toprule
    \textbf{Condition} & \textbf{Manifest change} & \textbf{Gate outcome} \\
    \midrule
    Baseline &
      permissive license \mbox{+} correct SHA-256 \mbox{+} audited source &
      \textbf{ADMITTED}: \texttt{license=Apache-2.0}, sha256 bound to bytes, source under
      audited prefix \\[4pt]

    Relabel license &
      \texttt{license = AGPL-3.0} &
      \textbf{REFUSED (license)}: \texttt{license "AGPL-3.0" is not allowed --- only
      permissive licenses accepted (Apache-2.0/MIT/BSD); AGPL is strictly forbidden} \\[4pt]

    Byte-swap weights &
      weights \(\neq\) declared SHA-256 &
      \textbf{REFUSED (hash)}: \texttt{sha256 mismatch for <model.onnx>: manifest declares
      aaaa\(\dots\)aaaa but file is 0b4895d0\(\dots\)29bb57 --- refusing to load (the
      declared license is bound to the audited bytes)} \\[4pt]

    Wrong origin &
      \texttt{source\_url} outside allowlist &
      \textbf{REFUSED (source)}: \texttt{source\_url
      "https://random-mirror.\allowbreak example.com/x.onnx" is not under any audited prefix in the
      verified-source allowlist} \\[4pt]

    Relabel vs.\ ledger &
      permissive but \texttt{license=MIT} where ledger says \texttt{Apache-2.0} &
      \textbf{REFUSED (ledger)}: \texttt{declared license "MIT" does not match the
      maintainer-audited upstream license "Apache-2.0" \dots\ (manifest may be mislabeled or
      relicensed)}; the allowlist cannot catch this (MIT is permissive), only the ledger does \\[4pt]

    Un-audited upstream &
      \texttt{source\_url} with no ledger entry &
      \textbf{REFUSED (ledger)}: \texttt{\dots\ has no maintainer-audited license-ledger entry
      (verified mode requires an audited upstream)} \\
    \bottomrule
  \end{tabular}
\end{table}

\subsection{Manifest Specification}
\label{sec:catalog:manifest}

Each model is described by a YAML \emph{manifest} file that the registry loads at
startup.  The manifest declares all information needed to validate the license,
locate ONNX artefacts, configure the execution-provider fallback chain, and pass
preprocessing parameters to the model's Go package without hard-coding them.

For plain (single-session) models, the manifest provides a single \texttt{model\_file}
path.  For \texttt{PipelineModel} instances that require multiple ONNX sessions
(e.g.\ MobileSAM encoder/decoder, Grounded-SAM), the manifest uses a \texttt{files:}
role map instead; each key is a logical role name that the model package uses to
request the correct session from \texttt{lifecycle.Manager}.

The following listing shows a representative manifest for \texttt{rf-detr-nano}:

\begin{lstlisting}[language=yaml, float=tbp, caption={Sample manifest: \texttt{rf-detr-nano}.},
                   label={lst:manifest}]
name: rf-detr-nano
task: detection
license: Apache-2.0
source_url: https://huggingface.co/PierreMarieCurie/.../rf-detr-nano.onnx
sha256: 3fcbba0f...be67abb2   # binds license to these exact bytes (verified mode)

execution_providers:
  - tensorrt
  - cuda
  - cpu

model_file: models/rf_detr_nano.onnx

preprocessing:
  resize: [560, 560]
  mean:   [0.485, 0.456, 0.406]
  std:    [0.229, 0.224, 0.225]
  input_format: NCHW
\end{lstlisting}

\noindent
The \texttt{execution\_providers} list is tried in order; if the ORT build on the host
does not include a listed provider, it is skipped silently and the next entry is
attempted.  \texttt{cpu} is always appended last to guarantee a functional fallback on
any hardware.  The \texttt{preprocessing} block is parsed into a \texttt{PreprocessMeta}
struct and passed to the model's \texttt{Preprocess} method, ensuring that bounding-box
coordinates are mapped back to original-image space before the result is serialised.

% No \begin{document} / \end{document}

\subsection{Task Types}
\label{sec:catalog:tasks}

VisionServe organises its model portfolio around seven computer-vision task types.
Table~\ref{tab:task-types} defines each type, states the expected input and output
contract, and names representative models currently registered in the system.

\begin{table*}[t]
  \centering
  \small
  \caption{Task types supported by VisionServe.}
  \label{tab:task-types}
  \begin{tabular}{llp{6.2cm}l}
    \toprule
    \textbf{Task} & \textbf{I/O Contract} & \textbf{Description} & \textbf{Example Models} \\
    \midrule
    \texttt{detection}     & image $\to$ boxes     &
      Localise and classify objects in an image.
      Returns axis-aligned bounding boxes in original-image coordinates
      together with class labels and confidence scores.
      &
      RF-DETR, SCRFD, PaddleOCR \\[4pt]

    \texttt{segmentation}  & image + prompt $\to$ masks &
      Produce per-pixel binary masks, optionally conditioned on a box or
      point prompt.  Masks are encoded as column-major run-length encoding
      (RLE) in the unified \texttt{Result} schema.
      &
      MobileSAM, EfficientSAM, SAM2 \\[4pt]

    \texttt{open\_vocab}   & image + text $\to$ boxes &
      Detect objects whose categories are specified as free-form natural-language
      queries at request time, without any fixed class vocabulary.
      &
      GroundingDINO \\[4pt]

    \texttt{classification}& image $\to$ label + score &
      Assign a single top-\(k\) label from a closed-set vocabulary
      (e.g.\ ImageNet-1k) to the entire image.
      &
      EfficientNet-B0, MobileNetV3 \\[4pt]

    \texttt{depth}         & image $\to$ depth map &
      Estimate a per-pixel relative or metric depth map.  Output is a
      single-channel float32 tensor aligned to the input resolution.
      &
      MiDaS, Depth Anything V2 \\[4pt]

    \texttt{embed}         & image $\to$ vector &
      Extract a compact, semantically rich embedding vector suitable for
      retrieval, zero-shot classification, or downstream ML pipelines.
      &
      CLIP ViT-B/32 \\[4pt]

    \texttt{grasp}         & image $\to$ grasps &
      Synthesise planar parallel-jaw grasps ($(x,y,\theta,\text{width},\text{quality})$ in
      original-image coordinates) from a segmentation mask via a pure-Go analytic
      antipodal force-closure search (no extra ONNX weights).  Class-agnostic
      (whole-image automask) or class-aware (detector $\to$ mask $\to$ grasp).
      &
      grasp, grasp-rfdetr, grasp-gd \\
    \bottomrule
  \end{tabular}
\end{table*}

\subsection{Model Inventory}
\label{sec:catalog:inventory}

Table~\ref{tab:model-inventory} enumerates all 19 models in the current VisionServe
registry.  Each model is distributed as one or more ONNX files; multi-session models
(MobileSAM, EfficientSAM, SAM2, Grounded-SAM, and the grasp pipelines) use the
\texttt{PipelineModel} interface and declare a \texttt{files:} role map in their manifest.
The three grasp pipelines reuse the segmentation/detection weights of sibling models and add
no new ONNX file: the grasp synthesis itself is a pure-Go analytic stage.

\begin{table*}[t]
  \centering
  \small
  \caption{%
    VisionServe model inventory.
    Architecture notes describe the primary network family.
    Models marked $\dagger$ are pending benchmark runs (see text).%
  }
  \label{tab:model-inventory}
  \begin{tabular}{llll}
    \toprule
    \textbf{Model} & \textbf{Task} & \textbf{License} & \textbf{Architecture} \\
    \midrule
    \texttt{rf-detr}          & detection     & Apache-2.0 & ResNet-50 backbone + DETR head (NMS-free) \\
    \texttt{rf-detr-nano}     & detection     & Apache-2.0 & Lightweight DETR variant \\
    \texttt{rt-detr}$^\dagger$& detection     & Apache-2.0 & Hybrid CNN/Transformer real-time DETR \\
    \texttt{scrfd}            & detection     & MIT        & Sample- \& computation-redistributed face detector \\
    \texttt{paddle-ocr}       & detection     & Apache-2.0 & PP-OCRv4 text detection backbone \\
    \texttt{grounding-dino}   & open\_vocab   & Apache-2.0 & Swin-T image encoder + BERT text encoder \\
    \texttt{mobile-sam}       & segmentation  & Apache-2.0 & Tiny ViT encoder + SAM decoder \\
    \texttt{efficient-sam}    & segmentation  & Apache-2.0 & ViT-S/16 encoder with masked-image pretraining \\
    \texttt{sam2}             & segmentation  & Apache-2.0 & Hiera-Tiny image/video encoder + SAM2 decoder \\
    \texttt{nano-sam}$^\dagger$& segmentation & Apache-2.0 & ResNet-18 distilled SAM encoder \\
    \texttt{grounded-sam}    & open\_vocab    & Apache-2.0 & GroundingDINO boxes $\to$ MobileSAM masks (composite) \\
    \texttt{midas}            & depth         & MIT        & DPT-Hybrid (ViT-H + dense prediction transformer) \\
    \texttt{depth-anything-v2}$^\dagger$ & depth & Apache-2.0 & ViT-S teacher-student depth foundation model \\
    \texttt{efficientnet-b0}  & classification& Apache-2.0 & EfficientNet-B0 compound-scaled CNN \\
    \texttt{mobilenet-v3}     & classification& Apache-2.0 & MobileNetV3-Large depthwise-separable CNN \\
    \texttt{clip}             & embed         & MIT        & CLIP ViT-B/32 vision encoder (image side only) \\
    \texttt{grasp}            & grasp         & Apache-2.0 & MobileSAM automask $\to$ pure-Go analytic grasps (class-agnostic) \\
    \texttt{grasp-rfdetr}     & grasp         & Apache-2.0 & RF-DETR boxes $\to$ MobileSAM $\to$ analytic grasps (class-aware) \\
    \texttt{grasp-gd}         & grasp         & Apache-2.0 & GroundingDINO boxes $\to$ MobileSAM $\to$ analytic grasps (text-prompted) \\
    \bottomrule
  \end{tabular}
\end{table*}

\noindent
Three models marked $\dagger$ remain pending: \texttt{rt-detr} and
\texttt{depth-anything-v2} require authenticated HuggingFace downloads, and
\texttt{nano-sam} requires a manual ONNX export.  All four grasp/composite pipelines
(\texttt{grasp}, \texttt{grasp-rfdetr}, \texttt{grounded-sam}, \texttt{grasp-gd}) are now
benchmarked (Section~\ref{sec:eval:latency}); the two GroundingDINO-based pipelines are
serialized for concurrency-safety (a fixed \texttt{SA\_ONSTACK} crash, Section~\ref{sec:limitations})
and are reported single-stream.

\paragraph{Grasp and license discipline.}
The grasp capability also illustrates why the license gate (C1) is load-bearing rather than
decorative.  In selecting a grasp method we rejected the most accurate learned grasp
models, GQ-CNN / FC-GQ-CNN (Dex-Net), under a UC~Berkeley ``educational, research, and
not-for-profit'' license, because that non-commercial bar is incompatible with
VisionServe's permissive, commercial-friendly stance, as AGPL is.  The shipped path
instead pairs a permissive MobileSAM mask with a clean-room pure-Go analytic synthesizer (no
weights, no license at all), so the seventh task adds capability without adding license risk.

\subsection{License Safety}
\label{sec:catalog:license}

Open-source computer-vision models carry heterogeneous licenses, and the distinction
between permissive and copyleft licenses has material consequences for downstream
deployments.  The most prominent example is the Ultralytics YOLO family
\cite{ultralytics2023yolo}, which is released under the GNU Affero General
Public License v3.0 (AGPL-3.0) \cite{agpllicense}.  AGPL-3.0 extends the GPL copyleft
to networked services: any software that uses AGPL code and provides access to
users over a network must release its entire source under AGPL-3.0.  As a result,
integrating a single AGPL model into an inference server would virally relicense the
server (and every product built on it) under AGPL-3.0.  This is incompatible with
closed-source, proprietary, or commercial deployments; the only escape is a separate
commercial license negotiated with the copyright holder.

VisionServe enforces an allowlist of permissive licenses
(\texttt{Apache-2.0}, \texttt{MIT}, \texttt{BSD-2-Clause}, \texttt{BSD-3-Clause})
in the registry parser at process startup \cite{apachelicense}.  Any model manifest
that declares a \texttt{license} field outside this allowlist causes the server to
abort with a clear error message, preventing accidental integration of copyleft models.
This check is implemented in code, not merely documented: the manifest
\texttt{license} field is a required string, and the parser rejects absent or
non-allowlisted values before any ONNX session is initialised.  This design choice is
deliberate: because VisionServe is itself Apache-2.0 and targets edge deployments in
commercial and closed-source products, its permissive character must be preserved by
every model it ships.

\paragraph{Threat model.}
We scope the license gate (contribution~C1), since its value depends on
stating both what it defends against and what it cannot.
\emph{In scope.}  The gate is designed to prevent four failure modes at the moment the
runtime loads model weights:
(i)~accidentally loading a copyleft (AGPL) model that would virally relicense an otherwise
permissive deployment;
(ii)~``permissive-washing'' or license drift, where a model's restrictive license is
silently dropped or relabelled as permissive on the way into the deployment;
(iii)~tampered or swapped weight bytes that no longer correspond to the audited model
(a content mismatch);
(iv)~weights pulled from an unvetted origin; and
(v)~a manifest relabelled by its author with a permissive-but-wrong license that diverges from
the maintainer-audited upstream (caught by the provenance ledger, which the allowlist alone
cannot detect because the wrong label is itself permissive).
\emph{Mechanism.}  Enforcement is entirely offline, in a single binary, and consists of
four bound checks applied at load time.
First, the declared \texttt{license} field must be in the permissive allowlist
(\texttt{Apache-2.0}/\texttt{MIT}/\texttt{BSD}); the registry parser
(\texttt{Manifest.validate()}) refuses any other value before a session is created.
Second, the SHA-256 digest of the weight bytes must equal the digest declared in the
manifest (\texttt{Manifest.VerifyWeights()}); a mismatch refuses the load.
Third, when a verified-source allowlist is supplied, the manifest's \texttt{source\_url}
must fall under an audited prefix.
Fourth, the check that closes the upstream-relabelling hole for the curated
catalog: in the opt-in verified mode the contributor-declared license must equal the
license recorded for that \texttt{source\_url} prefix in a \emph{maintainer-audited
license-provenance ledger} (\texttt{Manifest.VerifyLicenseProvenance()}).  The ledger is a
separate in-binary table in which a maintainer records, per audited upstream, the SPDX license
they verified by reading the upstream's actual LICENSE file, the URL of that file, and its
SHA-256 as evidence.  Because contributors edit manifests and maintainers edit the ledger, the
two authorities are separated: a pull request that flips a manifest's \texttt{license} to
sneak in a copyleft model is refused by a record the PR author does not control.  Together these
checks bind a declared-permissive license to specific audited bytes, an audited origin,
and an independent audited record of the upstream's true license, without any network call.
\emph{Out of scope.}  The allowlist and SHA-256/source pins alone still trust the declared
license string and mitigate only substitution (swapped bytes, wrong mirror), not
upstream mislabelling; the ledger removes that residue for the shipped catalog by
cross-checking against an independent human audit.  What remains out of scope is genuine: the
ledger's ground truth is a one-time human audit of the upstream LICENSE file (a re-audit detects
later divergence via the pinned LICENSE-file hash, but an upstream mislabelled at audit
time is not caught), and it covers only audited upstreams, not arbitrary un-audited uploads, for
which a universal license oracle is impossible.  We therefore frame~C1 as load-time
license-policy enforcement, hardened by content-hashing, a verified-source pin, and an
audited provenance ledger: strong for the curated catalog, and not a
machine-checkable license guarantee for arbitrary models.
\emph{Why existing layers do not cover it.}  As discussed in
Section~\ref{sec:related}, build-time software-composition-analysis scanners
(ScanCode, FOSSology, Snyk) operate on source dependency trees and emit reports, not a
load-time decision over weights; deploy-time admission controllers (Kyverno, OPA
Gatekeeper, Sigstore policy-controller) verify container-image signatures at pod
admission, i.e.\ container integrity, not a license policy over model weights; and model
signing (OpenSSF/Sigstore) proves signer identity and byte integrity, not license policy.
None of these enforces a license policy at the moment the runtime loads model weights,
offline, which is the gap~C1 fills.

\paragraph{Adversarial demonstration.}
To show that the gate refuses bad models rather than merely documenting an
intent to, Table~\ref{tab:gate} reports a reproducible adversarial demonstration.
Starting from a baseline manifest (permissive license, correct SHA-256, audited source)
that is admitted, we apply one adversarial mutation at a time and record the gate's
verdict.  Each mutation maps to one of the four bound checks above, and each is refused
with the runtime's own diagnostic message.  The six cases are exercised by the test
\texttt{TestGate\_AdversarialDemo}; the transcript is reproducible via
\texttt{go test ./internal/registry -run TestGate\_AdversarialDemo -v}.

\begin{table}[t]
  \centering
  \footnotesize
  \setlength{\tabcolsep}{3pt}
  \caption{%
    Reproducible adversarial demonstration of the load-time license gate (C1).
    Starting from an admitted baseline, each row applies one manifest mutation and
    records the runtime's verdict and (abbreviated) diagnostic.  Reproducible via
    \texttt{go test ./internal/registry -run TestGate\_AdversarialDemo -v}
    (\texttt{internal/registry/gate\_demo\_test.go}); the long digest is abbreviated
    with \(\dots\) and the temporary weight path is shown as \texttt{<model.onnx>}.%
  }
  \label{tab:gate}
  \begin{tabular}{p{1.2cm}p{1.5cm}p{5.4cm}}
    \toprule
    \textbf{Condition} & \textbf{Manifest change} & \textbf{Gate outcome} \\
    \midrule
    Baseline &
      permissive license \mbox{+} correct SHA-256 \mbox{+} audited source &
      \textbf{ADMITTED}: \texttt{license=Apache-2.0}, sha256 bound to bytes, source under
      audited prefix \\[4pt]

    Relabel license &
      \texttt{license = AGPL-3.0} &
      \textbf{REFUSED (license)}: \texttt{license "AGPL-3.0" is not allowed --- only
      permissive licenses accepted (Apache-2.0/MIT/BSD); AGPL is strictly forbidden} \\[4pt]

    Byte-swap weights &
      weights \(\neq\) declared SHA-256 &
      \textbf{REFUSED (hash)}: \texttt{sha256 mismatch for <model.onnx>: manifest declares
      aaaa\(\dots\)aaaa but file is 0b4895d0\(\dots\)29bb57 --- refusing to load (the
      declared license is bound to the audited bytes)} \\[4pt]

    Wrong origin &
      \texttt{source\_url} outside allowlist &
      \textbf{REFUSED (source)}: \texttt{source\_url
      "https://random-mirror.\allowbreak example.com/x.onnx" is not under any audited prefix in the
      verified-source allowlist} \\[4pt]

    Relabel vs.\ ledger &
      permissive but \texttt{license=MIT} where ledger says \texttt{Apache-2.0} &
      \textbf{REFUSED (ledger)}: \texttt{declared license "MIT" does not match the
      maintainer-audited upstream license "Apache-2.0" \dots\ (manifest may be mislabeled or
      relicensed)}; the allowlist cannot catch this (MIT is permissive), only the ledger does \\[4pt]

    Un-audited upstream &
      \texttt{source\_url} with no ledger entry &
      \textbf{REFUSED (ledger)}: \texttt{\dots\ has no maintainer-audited license-ledger entry
      (verified mode requires an audited upstream)} \\
    \bottomrule
  \end{tabular}
\end{table}

\subsection{Manifest Specification}
\label{sec:catalog:manifest}

Each model is described by a YAML \emph{manifest} file that the registry loads at
startup.  The manifest declares all information needed to validate the license,
locate ONNX artefacts, configure the execution-provider fallback chain, and pass
preprocessing parameters to the model's Go package without hard-coding them.

For plain (single-session) models, the manifest provides a single \texttt{model\_file}
path.  For \texttt{PipelineModel} instances that require multiple ONNX sessions
(e.g.\ MobileSAM encoder/decoder, Grounded-SAM), the manifest uses a \texttt{files:}
role map instead; each key is a logical role name that the model package uses to
request the correct session from \texttt{lifecycle.Manager}.

The following listing shows a representative manifest for \texttt{rf-detr-nano}:

\begin{lstlisting}[language=yaml, float=tbp, caption={Sample manifest: \texttt{rf-detr-nano}.},
                   label={lst:manifest}]
name: rf-detr-nano
task: detection
license: Apache-2.0
source_url: https://huggingface.co/PierreMarieCurie/.../rf-detr-nano.onnx
sha256: 3fcbba0f...be67abb2   # binds license to these exact bytes (verified mode)

execution_providers:
  - tensorrt
  - cuda
  - cpu

model_file: models/rf_detr_nano.onnx

preprocessing:
  resize: [560, 560]
  mean:   [0.485, 0.456, 0.406]
  std:    [0.229, 0.224, 0.225]
  input_format: NCHW
\end{lstlisting}

\noindent
The \texttt{execution\_providers} list is tried in order; if the ORT build on the host
does not include a listed provider, it is skipped silently and the next entry is
attempted.  \texttt{cpu} is always appended last to guarantee a functional fallback on
any hardware.  The \texttt{preprocessing} block is parsed into a \texttt{PreprocessMeta}
struct and passed to the model's \texttt{Preprocess} method, ensuring that bounding-box
coordinates are mapped back to original-image space before the result is serialised.

% No \begin{document} / \end{document}

\subsection{Task Types}
\label{sec:catalog:tasks}

VisionServe organises its model portfolio around seven computer-vision task types.
Table~\ref{tab:task-types} defines each type, states the expected input and output
contract, and names representative models currently registered in the system.

\begin{table*}[t]
  \centering
  \small
  \caption{Task types supported by VisionServe.}
  \label{tab:task-types}
  \begin{tabular}{llp{6.2cm}l}
    \toprule
    \textbf{Task} & \textbf{I/O Contract} & \textbf{Description} & \textbf{Example Models} \\
    \midrule
    \texttt{detection}     & image $\to$ boxes     &
      Localise and classify objects in an image.
      Returns axis-aligned bounding boxes in original-image coordinates
      together with class labels and confidence scores.
      &
      RF-DETR, SCRFD, PaddleOCR \\[4pt]

    \texttt{segmentation}  & image + prompt $\to$ masks &
      Produce per-pixel binary masks, optionally conditioned on a box or
      point prompt.  Masks are encoded as column-major run-length encoding
      (RLE) in the unified \texttt{Result} schema.
      &
      MobileSAM, EfficientSAM, SAM2 \\[4pt]

    \texttt{open\_vocab}   & image + text $\to$ boxes &
      Detect objects whose categories are specified as free-form natural-language
      queries at request time, without any fixed class vocabulary.
      &
      GroundingDINO \\[4pt]

    \texttt{classification}& image $\to$ label + score &
      Assign a single top-\(k\) label from a closed-set vocabulary
      (e.g.\ ImageNet-1k) to the entire image.
      &
      EfficientNet-B0, MobileNetV3 \\[4pt]

    \texttt{depth}         & image $\to$ depth map &
      Estimate a per-pixel relative or metric depth map.  Output is a
      single-channel float32 tensor aligned to the input resolution.
      &
      MiDaS, Depth Anything V2 \\[4pt]

    \texttt{embed}         & image $\to$ vector &
      Extract a compact, semantically rich embedding vector suitable for
      retrieval, zero-shot classification, or downstream ML pipelines.
      &
      CLIP ViT-B/32 \\[4pt]

    \texttt{grasp}         & image $\to$ grasps &
      Synthesise planar parallel-jaw grasps ($(x,y,\theta,\text{width},\text{quality})$ in
      original-image coordinates) from a segmentation mask via a pure-Go analytic
      antipodal force-closure search (no extra ONNX weights).  Class-agnostic
      (whole-image automask) or class-aware (detector $\to$ mask $\to$ grasp).
      &
      grasp, grasp-rfdetr, grasp-gd \\
    \bottomrule
  \end{tabular}
\end{table*}

\subsection{Model Inventory}
\label{sec:catalog:inventory}

Table~\ref{tab:model-inventory} enumerates all 19 models in the current VisionServe
registry.  Each model is distributed as one or more ONNX files; multi-session models
(MobileSAM, EfficientSAM, SAM2, Grounded-SAM, and the grasp pipelines) use the
\texttt{PipelineModel} interface and declare a \texttt{files:} role map in their manifest.
The three grasp pipelines reuse the segmentation/detection weights of sibling models and add
no new ONNX file: the grasp synthesis itself is a pure-Go analytic stage.

\begin{table*}[t]
  \centering
  \small
  \caption{%
    VisionServe model inventory.
    Architecture notes describe the primary network family.
    Models marked $\dagger$ are pending benchmark runs (see text).%
  }
  \label{tab:model-inventory}
  \begin{tabular}{llll}
    \toprule
    \textbf{Model} & \textbf{Task} & \textbf{License} & \textbf{Architecture} \\
    \midrule
    \texttt{rf-detr}          & detection     & Apache-2.0 & ResNet-50 backbone + DETR head (NMS-free) \\
    \texttt{rf-detr-nano}     & detection     & Apache-2.0 & Lightweight DETR variant \\
    \texttt{rt-detr}$^\dagger$& detection     & Apache-2.0 & Hybrid CNN/Transformer real-time DETR \\
    \texttt{scrfd}            & detection     & MIT        & Sample- \& computation-redistributed face detector \\
    \texttt{paddle-ocr}       & detection     & Apache-2.0 & PP-OCRv4 text detection backbone \\
    \texttt{grounding-dino}   & open\_vocab   & Apache-2.0 & Swin-T image encoder + BERT text encoder \\
    \texttt{mobile-sam}       & segmentation  & Apache-2.0 & Tiny ViT encoder + SAM decoder \\
    \texttt{efficient-sam}    & segmentation  & Apache-2.0 & ViT-S/16 encoder with masked-image pretraining \\
    \texttt{sam2}             & segmentation  & Apache-2.0 & Hiera-Tiny image/video encoder + SAM2 decoder \\
    \texttt{nano-sam}$^\dagger$& segmentation & Apache-2.0 & ResNet-18 distilled SAM encoder \\
    \texttt{grounded-sam}    & open\_vocab    & Apache-2.0 & GroundingDINO boxes $\to$ MobileSAM masks (composite) \\
    \texttt{midas}            & depth         & MIT        & DPT-Hybrid (ViT-H + dense prediction transformer) \\
    \texttt{depth-anything-v2}$^\dagger$ & depth & Apache-2.0 & ViT-S teacher-student depth foundation model \\
    \texttt{efficientnet-b0}  & classification& Apache-2.0 & EfficientNet-B0 compound-scaled CNN \\
    \texttt{mobilenet-v3}     & classification& Apache-2.0 & MobileNetV3-Large depthwise-separable CNN \\
    \texttt{clip}             & embed         & MIT        & CLIP ViT-B/32 vision encoder (image side only) \\
    \texttt{grasp}            & grasp         & Apache-2.0 & MobileSAM automask $\to$ pure-Go analytic grasps (class-agnostic) \\
    \texttt{grasp-rfdetr}     & grasp         & Apache-2.0 & RF-DETR boxes $\to$ MobileSAM $\to$ analytic grasps (class-aware) \\
    \texttt{grasp-gd}         & grasp         & Apache-2.0 & GroundingDINO boxes $\to$ MobileSAM $\to$ analytic grasps (text-prompted) \\
    \bottomrule
  \end{tabular}
\end{table*}

\noindent
Three models marked $\dagger$ remain pending: \texttt{rt-detr} and
\texttt{depth-anything-v2} require authenticated HuggingFace downloads, and
\texttt{nano-sam} requires a manual ONNX export.  All four grasp/composite pipelines
(\texttt{grasp}, \texttt{grasp-rfdetr}, \texttt{grounded-sam}, \texttt{grasp-gd}) are now
benchmarked (Section~\ref{sec:eval:latency}); the two GroundingDINO-based pipelines are
serialized for concurrency-safety (a fixed \texttt{SA\_ONSTACK} crash, Section~\ref{sec:limitations})
and are reported single-stream.

\paragraph{Grasp and license discipline.}
The grasp capability also illustrates why the license gate (C1) is load-bearing rather than
decorative.  In selecting a grasp method we rejected the most accurate learned grasp
models, GQ-CNN / FC-GQ-CNN (Dex-Net), under a UC~Berkeley ``educational, research, and
not-for-profit'' license, because that non-commercial bar is incompatible with
VisionServe's permissive, commercial-friendly stance, as AGPL is.  The shipped path
instead pairs a permissive MobileSAM mask with a clean-room pure-Go analytic synthesizer (no
weights, no license at all), so the seventh task adds capability without adding license risk.

\subsection{License Safety}
\label{sec:catalog:license}

Open-source computer-vision models carry heterogeneous licenses, and the distinction
between permissive and copyleft licenses has material consequences for downstream
deployments.  The most prominent example is the Ultralytics YOLO family
\cite{ultralytics2023yolo}, which is released under the GNU Affero General
Public License v3.0 (AGPL-3.0) \cite{agpllicense}.  AGPL-3.0 extends the GPL copyleft
to networked services: any software that uses AGPL code and provides access to
users over a network must release its entire source under AGPL-3.0.  As a result,
integrating a single AGPL model into an inference server would virally relicense the
server (and every product built on it) under AGPL-3.0.  This is incompatible with
closed-source, proprietary, or commercial deployments; the only escape is a separate
commercial license negotiated with the copyright holder.

VisionServe enforces an allowlist of permissive licenses
(\texttt{Apache-2.0}, \texttt{MIT}, \texttt{BSD-2-Clause}, \texttt{BSD-3-Clause})
in the registry parser at process startup \cite{apachelicense}.  Any model manifest
that declares a \texttt{license} field outside this allowlist causes the server to
abort with a clear error message, preventing accidental integration of copyleft models.
This check is implemented in code, not merely documented: the manifest
\texttt{license} field is a required string, and the parser rejects absent or
non-allowlisted values before any ONNX session is initialised.  This design choice is
deliberate: because VisionServe is itself Apache-2.0 and targets edge deployments in
commercial and closed-source products, its permissive character must be preserved by
every model it ships.

\paragraph{Threat model.}
We scope the license gate (contribution~C1), since its value depends on
stating both what it defends against and what it cannot.
\emph{In scope.}  The gate is designed to prevent four failure modes at the moment the
runtime loads model weights:
(i)~accidentally loading a copyleft (AGPL) model that would virally relicense an otherwise
permissive deployment;
(ii)~``permissive-washing'' or license drift, where a model's restrictive license is
silently dropped or relabelled as permissive on the way into the deployment;
(iii)~tampered or swapped weight bytes that no longer correspond to the audited model
(a content mismatch);
(iv)~weights pulled from an unvetted origin; and
(v)~a manifest relabelled by its author with a permissive-but-wrong license that diverges from
the maintainer-audited upstream (caught by the provenance ledger, which the allowlist alone
cannot detect because the wrong label is itself permissive).
\emph{Mechanism.}  Enforcement is entirely offline, in a single binary, and consists of
four bound checks applied at load time.
First, the declared \texttt{license} field must be in the permissive allowlist
(\texttt{Apache-2.0}/\texttt{MIT}/\texttt{BSD}); the registry parser
(\texttt{Manifest.validate()}) refuses any other value before a session is created.
Second, the SHA-256 digest of the weight bytes must equal the digest declared in the
manifest (\texttt{Manifest.VerifyWeights()}); a mismatch refuses the load.
Third, when a verified-source allowlist is supplied, the manifest's \texttt{source\_url}
must fall under an audited prefix.
Fourth, the check that closes the upstream-relabelling hole for the curated
catalog: in the opt-in verified mode the contributor-declared license must equal the
license recorded for that \texttt{source\_url} prefix in a \emph{maintainer-audited
license-provenance ledger} (\texttt{Manifest.VerifyLicenseProvenance()}).  The ledger is a
separate in-binary table in which a maintainer records, per audited upstream, the SPDX license
they verified by reading the upstream's actual LICENSE file, the URL of that file, and its
SHA-256 as evidence.  Because contributors edit manifests and maintainers edit the ledger, the
two authorities are separated: a pull request that flips a manifest's \texttt{license} to
sneak in a copyleft model is refused by a record the PR author does not control.  Together these
checks bind a declared-permissive license to specific audited bytes, an audited origin,
and an independent audited record of the upstream's true license, without any network call.
\emph{Out of scope.}  The allowlist and SHA-256/source pins alone still trust the declared
license string and mitigate only substitution (swapped bytes, wrong mirror), not
upstream mislabelling; the ledger removes that residue for the shipped catalog by
cross-checking against an independent human audit.  What remains out of scope is genuine: the
ledger's ground truth is a one-time human audit of the upstream LICENSE file (a re-audit detects
later divergence via the pinned LICENSE-file hash, but an upstream mislabelled at audit
time is not caught), and it covers only audited upstreams, not arbitrary un-audited uploads, for
which a universal license oracle is impossible.  We therefore frame~C1 as load-time
license-policy enforcement, hardened by content-hashing, a verified-source pin, and an
audited provenance ledger: strong for the curated catalog, and not a
machine-checkable license guarantee for arbitrary models.
\emph{Why existing layers do not cover it.}  As discussed in
Section~\ref{sec:related}, build-time software-composition-analysis scanners
(ScanCode, FOSSology, Snyk) operate on source dependency trees and emit reports, not a
load-time decision over weights; deploy-time admission controllers (Kyverno, OPA
Gatekeeper, Sigstore policy-controller) verify container-image signatures at pod
admission, i.e.\ container integrity, not a license policy over model weights; and model
signing (OpenSSF/Sigstore) proves signer identity and byte integrity, not license policy.
None of these enforces a license policy at the moment the runtime loads model weights,
offline, which is the gap~C1 fills.

\paragraph{Adversarial demonstration.}
To show that the gate refuses bad models rather than merely documenting an
intent to, Table~\ref{tab:gate} reports a reproducible adversarial demonstration.
Starting from a baseline manifest (permissive license, correct SHA-256, audited source)
that is admitted, we apply one adversarial mutation at a time and record the gate's
verdict.  Each mutation maps to one of the four bound checks above, and each is refused
with the runtime's own diagnostic message.  The six cases are exercised by the test
\texttt{TestGate\_AdversarialDemo}; the transcript is reproducible via
\texttt{go test ./internal/registry -run TestGate\_AdversarialDemo -v}.

\begin{table}[t]
  \centering
  \footnotesize
  \setlength{\tabcolsep}{3pt}
  \caption{%
    Reproducible adversarial demonstration of the load-time license gate (C1).
    Starting from an admitted baseline, each row applies one manifest mutation and
    records the runtime's verdict and (abbreviated) diagnostic.  Reproducible via
    \texttt{go test ./internal/registry -run TestGate\_AdversarialDemo -v}
    (\texttt{internal/registry/gate\_demo\_test.go}); the long digest is abbreviated
    with \(\dots\) and the temporary weight path is shown as \texttt{<model.onnx>}.%
  }
  \label{tab:gate}
  \begin{tabular}{p{1.2cm}p{1.5cm}p{5.4cm}}
    \toprule
    \textbf{Condition} & \textbf{Manifest change} & \textbf{Gate outcome} \\
    \midrule
    Baseline &
      permissive license \mbox{+} correct SHA-256 \mbox{+} audited source &
      \textbf{ADMITTED}: \texttt{license=Apache-2.0}, sha256 bound to bytes, source under
      audited prefix \\[4pt]

    Relabel license &
      \texttt{license = AGPL-3.0} &
      \textbf{REFUSED (license)}: \texttt{license "AGPL-3.0" is not allowed --- only
      permissive licenses accepted (Apache-2.0/MIT/BSD); AGPL is strictly forbidden} \\[4pt]

    Byte-swap weights &
      weights \(\neq\) declared SHA-256 &
      \textbf{REFUSED (hash)}: \texttt{sha256 mismatch for <model.onnx>: manifest declares
      aaaa\(\dots\)aaaa but file is 0b4895d0\(\dots\)29bb57 --- refusing to load (the
      declared license is bound to the audited bytes)} \\[4pt]

    Wrong origin &
      \texttt{source\_url} outside allowlist &
      \textbf{REFUSED (source)}: \texttt{source\_url
      "https://random-mirror.\allowbreak example.com/x.onnx" is not under any audited prefix in the
      verified-source allowlist} \\[4pt]

    Relabel vs.\ ledger &
      permissive but \texttt{license=MIT} where ledger says \texttt{Apache-2.0} &
      \textbf{REFUSED (ledger)}: \texttt{declared license "MIT" does not match the
      maintainer-audited upstream license "Apache-2.0" \dots\ (manifest may be mislabeled or
      relicensed)}; the allowlist cannot catch this (MIT is permissive), only the ledger does \\[4pt]

    Un-audited upstream &
      \texttt{source\_url} with no ledger entry &
      \textbf{REFUSED (ledger)}: \texttt{\dots\ has no maintainer-audited license-ledger entry
      (verified mode requires an audited upstream)} \\
    \bottomrule
  \end{tabular}
\end{table}

\subsection{Manifest Specification}
\label{sec:catalog:manifest}

Each model is described by a YAML \emph{manifest} file that the registry loads at
startup.  The manifest declares all information needed to validate the license,
locate ONNX artefacts, configure the execution-provider fallback chain, and pass
preprocessing parameters to the model's Go package without hard-coding them.

For plain (single-session) models, the manifest provides a single \texttt{model\_file}
path.  For \texttt{PipelineModel} instances that require multiple ONNX sessions
(e.g.\ MobileSAM encoder/decoder, Grounded-SAM), the manifest uses a \texttt{files:}
role map instead; each key is a logical role name that the model package uses to
request the correct session from \texttt{lifecycle.Manager}.

The following listing shows a representative manifest for \texttt{rf-detr-nano}:

\begin{lstlisting}[language=yaml, float=tbp, caption={Sample manifest: \texttt{rf-detr-nano}.},
                   label={lst:manifest}]
name: rf-detr-nano
task: detection
license: Apache-2.0
source_url: https://huggingface.co/PierreMarieCurie/.../rf-detr-nano.onnx
sha256: 3fcbba0f...be67abb2   # binds license to these exact bytes (verified mode)

execution_providers:
  - tensorrt
  - cuda
  - cpu

model_file: models/rf_detr_nano.onnx

preprocessing:
  resize: [560, 560]
  mean:   [0.485, 0.456, 0.406]
  std:    [0.229, 0.224, 0.225]
  input_format: NCHW
\end{lstlisting}

\noindent
The \texttt{execution\_providers} list is tried in order; if the ORT build on the host
does not include a listed provider, it is skipped silently and the next entry is
attempted.  \texttt{cpu} is always appended last to guarantee a functional fallback on
any hardware.  The \texttt{preprocessing} block is parsed into a \texttt{PreprocessMeta}
struct and passed to the model's \texttt{Preprocess} method, ensuring that bounding-box
coordinates are mapped back to original-image space before the result is serialised.

% No \begin{document} / \end{document}

\subsection{Task Types}
\label{sec:catalog:tasks}

VisionServe organises its model portfolio around seven computer-vision task types.
Table~\ref{tab:task-types} defines each type, states the expected input and output
contract, and names representative models currently registered in the system.

\begin{table*}[t]
  \centering
  \small
  \caption{Task types supported by VisionServe.}
  \label{tab:task-types}
  \begin{tabular}{llp{6.2cm}l}
    \toprule
    \textbf{Task} & \textbf{I/O Contract} & \textbf{Description} & \textbf{Example Models} \\
    \midrule
    \texttt{detection}     & image $\to$ boxes     &
      Localise and classify objects in an image.
      Returns axis-aligned bounding boxes in original-image coordinates
      together with class labels and confidence scores.
      &
      RF-DETR, SCRFD, PaddleOCR \\[4pt]

    \texttt{segmentation}  & image + prompt $\to$ masks &
      Produce per-pixel binary masks, optionally conditioned on a box or
      point prompt.  Masks are encoded as column-major run-length encoding
      (RLE) in the unified \texttt{Result} schema.
      &
      MobileSAM, EfficientSAM, SAM2 \\[4pt]

    \texttt{open\_vocab}   & image + text $\to$ boxes &
      Detect objects whose categories are specified as free-form natural-language
      queries at request time, without any fixed class vocabulary.
      &
      GroundingDINO \\[4pt]

    \texttt{classification}& image $\to$ label + score &
      Assign a single top-\(k\) label from a closed-set vocabulary
      (e.g.\ ImageNet-1k) to the entire image.
      &
      EfficientNet-B0, MobileNetV3 \\[4pt]

    \texttt{depth}         & image $\to$ depth map &
      Estimate a per-pixel relative or metric depth map.  Output is a
      single-channel float32 tensor aligned to the input resolution.
      &
      MiDaS, Depth Anything V2 \\[4pt]

    \texttt{embed}         & image $\to$ vector &
      Extract a compact, semantically rich embedding vector suitable for
      retrieval, zero-shot classification, or downstream ML pipelines.
      &
      CLIP ViT-B/32 \\[4pt]

    \texttt{grasp}         & image $\to$ grasps &
      Synthesise planar parallel-jaw grasps ($(x,y,\theta,\text{width},\text{quality})$ in
      original-image coordinates) from a segmentation mask via a pure-Go analytic
      antipodal force-closure search (no extra ONNX weights).  Class-agnostic
      (whole-image automask) or class-aware (detector $\to$ mask $\to$ grasp).
      &
      grasp, grasp-rfdetr, grasp-gd \\
    \bottomrule
  \end{tabular}
\end{table*}

\subsection{Model Inventory}
\label{sec:catalog:inventory}

Table~\ref{tab:model-inventory} enumerates all 19 models in the current VisionServe
registry.  Each model is distributed as one or more ONNX files; multi-session models
(MobileSAM, EfficientSAM, SAM2, Grounded-SAM, and the grasp pipelines) use the
\texttt{PipelineModel} interface and declare a \texttt{files:} role map in their manifest.
The three grasp pipelines reuse the segmentation/detection weights of sibling models and add
no new ONNX file: the grasp synthesis itself is a pure-Go analytic stage.

\begin{table*}[t]
  \centering
  \small
  \caption{%
    VisionServe model inventory.
    Architecture notes describe the primary network family.
    Models marked $\dagger$ are pending benchmark runs (see text).%
  }
  \label{tab:model-inventory}
  \begin{tabular}{llll}
    \toprule
    \textbf{Model} & \textbf{Task} & \textbf{License} & \textbf{Architecture} \\
    \midrule
    \texttt{rf-detr}          & detection     & Apache-2.0 & ResNet-50 backbone + DETR head (NMS-free) \\
    \texttt{rf-detr-nano}     & detection     & Apache-2.0 & Lightweight DETR variant \\
    \texttt{rt-detr}$^\dagger$& detection     & Apache-2.0 & Hybrid CNN/Transformer real-time DETR \\
    \texttt{scrfd}            & detection     & MIT        & Sample- \& computation-redistributed face detector \\
    \texttt{paddle-ocr}       & detection     & Apache-2.0 & PP-OCRv4 text detection backbone \\
    \texttt{grounding-dino}   & open\_vocab   & Apache-2.0 & Swin-T image encoder + BERT text encoder \\
    \texttt{mobile-sam}       & segmentation  & Apache-2.0 & Tiny ViT encoder + SAM decoder \\
    \texttt{efficient-sam}    & segmentation  & Apache-2.0 & ViT-S/16 encoder with masked-image pretraining \\
    \texttt{sam2}             & segmentation  & Apache-2.0 & Hiera-Tiny image/video encoder + SAM2 decoder \\
    \texttt{nano-sam}$^\dagger$& segmentation & Apache-2.0 & ResNet-18 distilled SAM encoder \\
    \texttt{grounded-sam}    & open\_vocab    & Apache-2.0 & GroundingDINO boxes $\to$ MobileSAM masks (composite) \\
    \texttt{midas}            & depth         & MIT        & DPT-Hybrid (ViT-H + dense prediction transformer) \\
    \texttt{depth-anything-v2}$^\dagger$ & depth & Apache-2.0 & ViT-S teacher-student depth foundation model \\
    \texttt{efficientnet-b0}  & classification& Apache-2.0 & EfficientNet-B0 compound-scaled CNN \\
    \texttt{mobilenet-v3}     & classification& Apache-2.0 & MobileNetV3-Large depthwise-separable CNN \\
    \texttt{clip}             & embed         & MIT        & CLIP ViT-B/32 vision encoder (image side only) \\
    \texttt{grasp}            & grasp         & Apache-2.0 & MobileSAM automask $\to$ pure-Go analytic grasps (class-agnostic) \\
    \texttt{grasp-rfdetr}     & grasp         & Apache-2.0 & RF-DETR boxes $\to$ MobileSAM $\to$ analytic grasps (class-aware) \\
    \texttt{grasp-gd}         & grasp         & Apache-2.0 & GroundingDINO boxes $\to$ MobileSAM $\to$ analytic grasps (text-prompted) \\
    \bottomrule
  \end{tabular}
\end{table*}

\noindent
Three models marked $\dagger$ remain pending: \texttt{rt-detr} and
\texttt{depth-anything-v2} require authenticated HuggingFace downloads, and
\texttt{nano-sam} requires a manual ONNX export.  All four grasp/composite pipelines
(\texttt{grasp}, \texttt{grasp-rfdetr}, \texttt{grounded-sam}, \texttt{grasp-gd}) are now
benchmarked (Section~\ref{sec:eval:latency}); the two GroundingDINO-based pipelines are
serialized for concurrency-safety (a fixed \texttt{SA\_ONSTACK} crash, Section~\ref{sec:limitations})
and are reported single-stream.

\paragraph{Grasp and license discipline.}
The grasp capability also illustrates why the license gate (C1) is load-bearing rather than
decorative.  In selecting a grasp method we rejected the most accurate learned grasp
models, GQ-CNN / FC-GQ-CNN (Dex-Net), under a UC~Berkeley ``educational, research, and
not-for-profit'' license, because that non-commercial bar is incompatible with
VisionServe's permissive, commercial-friendly stance, as AGPL is.  The shipped path
instead pairs a permissive MobileSAM mask with a clean-room pure-Go analytic synthesizer (no
weights, no license at all), so the seventh task adds capability without adding license risk.

\subsection{License Safety}
\label{sec:catalog:license}

Open-source computer-vision models carry heterogeneous licenses, and the distinction
between permissive and copyleft licenses has material consequences for downstream
deployments.  The most prominent example is the Ultralytics YOLO family
\cite{ultralytics2023yolo}, which is released under the GNU Affero General
Public License v3.0 (AGPL-3.0) \cite{agpllicense}.  AGPL-3.0 extends the GPL copyleft
to networked services: any software that uses AGPL code and provides access to
users over a network must release its entire source under AGPL-3.0.  As a result,
integrating a single AGPL model into an inference server would virally relicense the
server (and every product built on it) under AGPL-3.0.  This is incompatible with
closed-source, proprietary, or commercial deployments; the only escape is a separate
commercial license negotiated with the copyright holder.

VisionServe enforces an allowlist of permissive licenses
(\texttt{Apache-2.0}, \texttt{MIT}, \texttt{BSD-2-Clause}, \texttt{BSD-3-Clause})
in the registry parser at process startup \cite{apachelicense}.  Any model manifest
that declares a \texttt{license} field outside this allowlist causes the server to
abort with a clear error message, preventing accidental integration of copyleft models.
This check is implemented in code, not merely documented: the manifest
\texttt{license} field is a required string, and the parser rejects absent or
non-allowlisted values before any ONNX session is initialised.  This design choice is
deliberate: because VisionServe is itself Apache-2.0 and targets edge deployments in
commercial and closed-source products, its permissive character must be preserved by
every model it ships.

\paragraph{Threat model.}
We scope the license gate (contribution~C1), since its value depends on
stating both what it defends against and what it cannot.
\emph{In scope.}  The gate is designed to prevent four failure modes at the moment the
runtime loads model weights:
(i)~accidentally loading a copyleft (AGPL) model that would virally relicense an otherwise
permissive deployment;
(ii)~``permissive-washing'' or license drift, where a model's restrictive license is
silently dropped or relabelled as permissive on the way into the deployment;
(iii)~tampered or swapped weight bytes that no longer correspond to the audited model
(a content mismatch);
(iv)~weights pulled from an unvetted origin; and
(v)~a manifest relabelled by its author with a permissive-but-wrong license that diverges from
the maintainer-audited upstream (caught by the provenance ledger, which the allowlist alone
cannot detect because the wrong label is itself permissive).
\emph{Mechanism.}  Enforcement is entirely offline, in a single binary, and consists of
four bound checks applied at load time.
First, the declared \texttt{license} field must be in the permissive allowlist
(\texttt{Apache-2.0}/\texttt{MIT}/\texttt{BSD}); the registry parser
(\texttt{Manifest.validate()}) refuses any other value before a session is created.
Second, the SHA-256 digest of the weight bytes must equal the digest declared in the
manifest (\texttt{Manifest.VerifyWeights()}); a mismatch refuses the load.
Third, when a verified-source allowlist is supplied, the manifest's \texttt{source\_url}
must fall under an audited prefix.
Fourth, the check that closes the upstream-relabelling hole for the curated
catalog: in the opt-in verified mode the contributor-declared license must equal the
license recorded for that \texttt{source\_url} prefix in a \emph{maintainer-audited
license-provenance ledger} (\texttt{Manifest.VerifyLicenseProvenance()}).  The ledger is a
separate in-binary table in which a maintainer records, per audited upstream, the SPDX license
they verified by reading the upstream's actual LICENSE file, the URL of that file, and its
SHA-256 as evidence.  Because contributors edit manifests and maintainers edit the ledger, the
two authorities are separated: a pull request that flips a manifest's \texttt{license} to
sneak in a copyleft model is refused by a record the PR author does not control.  Together these
checks bind a declared-permissive license to specific audited bytes, an audited origin,
and an independent audited record of the upstream's true license, without any network call.
\emph{Out of scope.}  The allowlist and SHA-256/source pins alone still trust the declared
license string and mitigate only substitution (swapped bytes, wrong mirror), not
upstream mislabelling; the ledger removes that residue for the shipped catalog by
cross-checking against an independent human audit.  What remains out of scope is genuine: the
ledger's ground truth is a one-time human audit of the upstream LICENSE file (a re-audit detects
later divergence via the pinned LICENSE-file hash, but an upstream mislabelled at audit
time is not caught), and it covers only audited upstreams, not arbitrary un-audited uploads, for
which a universal license oracle is impossible.  We therefore frame~C1 as load-time
license-policy enforcement, hardened by content-hashing, a verified-source pin, and an
audited provenance ledger: strong for the curated catalog, and not a
machine-checkable license guarantee for arbitrary models.
\emph{Why existing layers do not cover it.}  As discussed in
Section~\ref{sec:related}, build-time software-composition-analysis scanners
(ScanCode, FOSSology, Snyk) operate on source dependency trees and emit reports, not a
load-time decision over weights; deploy-time admission controllers (Kyverno, OPA
Gatekeeper, Sigstore policy-controller) verify container-image signatures at pod
admission, i.e.\ container integrity, not a license policy over model weights; and model
signing (OpenSSF/Sigstore) proves signer identity and byte integrity, not license policy.
None of these enforces a license policy at the moment the runtime loads model weights,
offline, which is the gap~C1 fills.

\paragraph{Adversarial demonstration.}
To show that the gate refuses bad models rather than merely documenting an
intent to, Table~\ref{tab:gate} reports a reproducible adversarial demonstration.
Starting from a baseline manifest (permissive license, correct SHA-256, audited source)
that is admitted, we apply one adversarial mutation at a time and record the gate's
verdict.  Each mutation maps to one of the four bound checks above, and each is refused
with the runtime's own diagnostic message.  The six cases are exercised by the test
\texttt{TestGate\_AdversarialDemo}; the transcript is reproducible via
\texttt{go test ./internal/registry -run TestGate\_AdversarialDemo -v}.

\begin{table}[t]
  \centering
  \footnotesize
  \setlength{\tabcolsep}{3pt}
  \caption{%
    Reproducible adversarial demonstration of the load-time license gate (C1).
    Starting from an admitted baseline, each row applies one manifest mutation and
    records the runtime's verdict and (abbreviated) diagnostic.  Reproducible via
    \texttt{go test ./internal/registry -run TestGate\_AdversarialDemo -v}
    (\texttt{internal/registry/gate\_demo\_test.go}); the long digest is abbreviated
    with \(\dots\) and the temporary weight path is shown as \texttt{<model.onnx>}.%
  }
  \label{tab:gate}
  \begin{tabular}{p{1.2cm}p{1.5cm}p{5.4cm}}
    \toprule
    \textbf{Condition} & \textbf{Manifest change} & \textbf{Gate outcome} \\
    \midrule
    Baseline &
      permissive license \mbox{+} correct SHA-256 \mbox{+} audited source &
      \textbf{ADMITTED}: \texttt{license=Apache-2.0}, sha256 bound to bytes, source under
      audited prefix \\[4pt]

    Relabel license &
      \texttt{license = AGPL-3.0} &
      \textbf{REFUSED (license)}: \texttt{license "AGPL-3.0" is not allowed --- only
      permissive licenses accepted (Apache-2.0/MIT/BSD); AGPL is strictly forbidden} \\[4pt]

    Byte-swap weights &
      weights \(\neq\) declared SHA-256 &
      \textbf{REFUSED (hash)}: \texttt{sha256 mismatch for <model.onnx>: manifest declares
      aaaa\(\dots\)aaaa but file is 0b4895d0\(\dots\)29bb57 --- refusing to load (the
      declared license is bound to the audited bytes)} \\[4pt]

    Wrong origin &
      \texttt{source\_url} outside allowlist &
      \textbf{REFUSED (source)}: \texttt{source\_url
      "https://random-mirror.\allowbreak example.com/x.onnx" is not under any audited prefix in the
      verified-source allowlist} \\[4pt]

    Relabel vs.\ ledger &
      permissive but \texttt{license=MIT} where ledger says \texttt{Apache-2.0} &
      \textbf{REFUSED (ledger)}: \texttt{declared license "MIT" does not match the
      maintainer-audited upstream license "Apache-2.0" \dots\ (manifest may be mislabeled or
      relicensed)}; the allowlist cannot catch this (MIT is permissive), only the ledger does \\[4pt]

    Un-audited upstream &
      \texttt{source\_url} with no ledger entry &
      \textbf{REFUSED (ledger)}: \texttt{\dots\ has no maintainer-audited license-ledger entry
      (verified mode requires an audited upstream)} \\
    \bottomrule
  \end{tabular}
\end{table}

\subsection{Manifest Specification}
\label{sec:catalog:manifest}

Each model is described by a YAML \emph{manifest} file that the registry loads at
startup.  The manifest declares all information needed to validate the license,
locate ONNX artefacts, configure the execution-provider fallback chain, and pass
preprocessing parameters to the model's Go package without hard-coding them.

For plain (single-session) models, the manifest provides a single \texttt{model\_file}
path.  For \texttt{PipelineModel} instances that require multiple ONNX sessions
(e.g.\ MobileSAM encoder/decoder, Grounded-SAM), the manifest uses a \texttt{files:}
role map instead; each key is a logical role name that the model package uses to
request the correct session from \texttt{lifecycle.Manager}.

The following listing shows a representative manifest for \texttt{rf-detr-nano}:

\begin{lstlisting}[language=yaml, float=tbp, caption={Sample manifest: \texttt{rf-detr-nano}.},
                   label={lst:manifest}]
name: rf-detr-nano
task: detection
license: Apache-2.0
source_url: https://huggingface.co/PierreMarieCurie/.../rf-detr-nano.onnx
sha256: 3fcbba0f...be67abb2   # binds license to these exact bytes (verified mode)

execution_providers:
  - tensorrt
  - cuda
  - cpu

model_file: models/rf_detr_nano.onnx

preprocessing:
  resize: [560, 560]
  mean:   [0.485, 0.456, 0.406]
  std:    [0.229, 0.224, 0.225]
  input_format: NCHW
\end{lstlisting}

\noindent
The \texttt{execution\_providers} list is tried in order; if the ORT build on the host
does not include a listed provider, it is skipped silently and the next entry is
attempted.  \texttt{cpu} is always appended last to guarantee a functional fallback on
any hardware.  The \texttt{preprocessing} block is parsed into a \texttt{PreprocessMeta}
struct and passed to the model's \texttt{Preprocess} method, ensuring that bounding-box
coordinates are mapped back to original-image space before the result is serialised.
